\begin{abstract}
Advertisers usually target the users most likely to convert. The causal
question is different: which users convert \emph{because of} the ad? We work
through the full statistical-ML workflow on 13,979,592 users from Criteo's
randomized incrementality tests (v2.1), from exploratory analysis and
hypothesis tests through predictive models, average-treatment-effect (ATE)
estimators, heterogeneous-effect (CATE) and uplift learners, and a
budget-constrained targeting optimizer. Three findings stand out. First,
randomization is imperfect in the pooled benchmark: treatment is weakly
predictable from the features (classifier two-sample AUC 0.509), and
covariate-adjusted estimators agree with each other while sitting 13--28\,\%
below the raw treated-versus-control difference (doubly robust visit effect
+0.746\,pp against a raw +1.034\,pp). Second, assignment is not exposure: only
3.6\,\% of eligible users were shown an ad, and using assignment as an
instrument the effect of actually seeing an ad is +21.2\,pp on visits and
+2.84\,pp on conversions. Third, the central experiment has a split answer.
For visits, uplift targeting at a 5\,\% budget finds 13,624 incremental
visits against 8,542 for ranking by visit probability (+60\,\%); for
conversions, no uplift learner beats ranking by conversion probability at any
budget, because the ad lifts conversion roughly in proportion to baseline
propensity. Targeting nevertheless pays: treating the top 10\,\% by
conversion probability captures about 89\,\% of the lift of treating
everyone. All estimates use a fixed train/validation/test protocol with
bootstrap confidence intervals, and all code and results are released.

\medskip\noindent\small\textbf{Code:} \url{https://github.com/pleyva2004/causalml-criteo}\\
\textbf{Data:} \url{https://www.kaggle.com/datasets/arashnic/uplift-modeling}
\end{abstract}
